\chapter{Hành Trình Khám Phá Bản Thân (The Journey of Self-Discovery)}
\begin{center}
  \textit{\color{truyengold}``Hành trình vạn dặm bắt đầu bằng một bước chân --- và bước chân đầu tiên phải đi vào trong lòng mình.''}\\[4pt]
  \textit{(A journey of a thousand miles begins with one step --- and the first step must go inward.)}\\[4pt]
  \small\color{truyenblue}--- Lão Tử ---
\end{center}
\ngancach

\section{Siddhartha Và Cây Bồ Đề}
\begin{truyen}{Siddhartha Và Cây Bồ Đề}{Phật Thích Ca Mâu Ni --- Ấn Độ, thế kỷ 5 TCN}
\chuhoa{H}{oàng} tử Siddhartha từ bỏ ngai vàng, vợ con, và cuộc sống xa hoa để đi tìm câu trả lời cho câu hỏi: tại sao con người phải khổ đau?\\[4pt]
\textit{(Prince Siddhartha abandoned his throne, family, and luxurious life to seek the answer to one question: why must humans suffer?)}

Sáu năm ông khổ hạnh tu luyện đến kiệt sức mà vẫn không tìm ra câu trả lời.\\[4pt]
\textit{(For six years he practiced extreme asceticism to exhaustion yet still found no answer.)}

Cuối cùng ông ngồi dưới cây bồ đề và nguyện không đứng dậy cho đến khi giác ngộ.\\[4pt]
\textit{(Finally he sat beneath the Bodhi tree and vowed not to rise until he attained enlightenment.)}

Câu trả lời ông tìm thấy không ở ngoài thế giới --- mà ở bên trong chính tâm mình: khổ đau đến từ chấp thủ, và giải thoát đến từ buông bỏ.\\[4pt]
\textit{(The answer he found was not in the outer world --- but within his own mind: suffering comes from clinging, and liberation comes from releasing.)}

\end{truyen}
\begin{baihoc}
  Những câu hỏi quan trọng nhất của cuộc đời thường không tìm được câu trả lời từ bên ngoài --- mà từ bên trong sự im lặng của chính mình.\\[4pt]
  \textit{(Life's most important questions often find no answer from the outside --- but from within one's own silence.)}
\end{baihoc}

\section{Marcus Aurelius Và Nhật Ký Riêng Tư}
\begin{truyen}{Marcus Aurelius Và Nhật Ký Riêng Tư}{Marcus Aurelius --- Đế Quốc La Mã, thế kỷ 2}
\chuhoa{M}{arcus} Aurelius là hoàng đế hùng mạnh nhất thế giới thời bấy giờ --- có quyền năng sinh sát tuyệt đối --- nhưng mỗi tối ông vẫn viết nhật ký để tự kiểm điểm mình.\\[4pt]
\textit{(Marcus Aurelius was the most powerful emperor in the world of his time --- with absolute power of life and death --- yet each evening he still wrote in his diary to examine himself.)}

Cuốn nhật ký đó --- sau này được đặt tên là ``Những Suy Tư'' (\textit{Meditations}) --- không bao giờ được viết để xuất bản, chỉ để tự nhắc nhở chính hoàng đế.\\[4pt]
\textit{(That diary --- later named ``Meditations'' --- was never written for publication, only to remind the emperor himself.)}

Ông viết: ``Mỗi sáng khi thức dậy, hãy nói với chính mình: Hôm nay tôi sẽ gặp những kẻ tò mò, vô ơn, thô lỗ, xảo quyệt và ích kỷ. Nhưng tôi vẫn không thể để điều đó làm hại tôi.''\\[4pt]
\textit{(He wrote: `Each morning when you wake, say to yourself: Today I will meet the meddling, ungrateful, overbearing, deceitful, envious, and unsociable. But I cannot let that harm me.')}

\end{truyen}
\begin{baihoc}
  Quyền năng thật sự không nằm ở việc điều khiển người khác --- mà ở việc điều khiển được phản ứng của chính mình.\\[4pt]
  \textit{(True power lies not in controlling others --- but in controlling one's own responses.)}
\end{baihoc}

\section{Nguyễn Trãi Viết Bình Ngô Đại Cáo}
\begin{truyen}{Nguyễn Trãi Viết Bình Ngô Đại Cáo}{Nguyễn Trãi --- Việt Nam, 1428}
\chuhoa{S}{au} mười năm kháng chiến gian khổ, khi quân Minh bị đánh đuổi, Nguyễn Trãi được Lê Lợi giao viết bản tuyên ngôn lịch sử.\\[4pt]
\textit{(After ten years of arduous resistance, when the Ming forces were expelled, Nguyen Trai was entrusted by Le Loi to write the historic proclamation.)}

Ông mở đầu bằng chân lý vĩnh cửu: ``Việc nhân nghĩa cốt ở yên dân --- quân điếu phạt trước lo trừ bạo.''\\[4pt]
\textit{(He opened with eternal truth: `The essence of humanity and righteousness lies in bringing peace to the people --- the foremost concern in raising arms is to rid them of tyranny.')}

Rồi ông khẳng định bản sắc dân tộc bằng những câu văn hùng hồn: ``Như nước Đại Việt ta từ trước, vốn xưng nền văn hiến đã lâu...''\\[4pt]
\textit{(Then he affirmed national identity with resounding lines: `Our Dai Viet, as before, has long proclaimed a culture of civilization...')}

Nguyễn Trãi không chỉ viết một bài cáo --- ông xây dựng ý thức về một dân tộc có bản sắc, có văn hóa, có chủ quyền riêng biệt.\\[4pt]
\textit{(Nguyen Trai did not merely write a proclamation --- he built consciousness of a people with distinct identity, culture, and sovereignty.)}

\end{truyen}
\begin{baihoc}
  Hiểu rõ mình là ai --- dân tộc mình là gì --- là nền tảng không gì lay chuyển được của lòng tự tôn và sức mạnh trường tồn.\\[4pt]
  \textit{(Clearly knowing who you are --- what your people stand for --- is the unshakeable foundation of self-respect and enduring strength.)}
\end{baihoc}

\section{Viktor Frankl Tìm Ý Nghĩa Trong Địa Ngục}
\begin{truyen}{Viktor Frankl Tìm Ý Nghĩa Trong Địa Ngục}{Viktor Frankl --- Áo, 1942--1945}
\chuhoa{V}{iktor} Frankl là bác sĩ tâm lý người Áo gốc Do Thái, bị đưa vào trại tập trung Auschwitz cùng toàn bộ gia đình; cha, mẹ, anh và vợ ông đều chết trong các trại.\\[4pt]
\textit{(Viktor Frankl was an Austrian Jewish psychiatrist, sent to Auschwitz along with his entire family; his father, mother, brother, and wife all perished in the camps.)}

Trong địa ngục đó, ông quan sát: những người sống sót không phải là những người khỏe nhất hay may mắn nhất --- mà là những người tìm được lý do để sống.\\[4pt]
\textit{(In that hell, he observed: the survivors were not the strongest or luckiest --- but those who found a reason to live.)}

Ông viết trong cuốn sách bất hủ \textit{Đi Tìm Ý Nghĩa Cuộc Sống}: ``Mọi thứ đều có thể bị lấy đi từ con người --- trừ một thứ: tự do cuối cùng của con người là lựa chọn thái độ của mình trước bất kỳ hoàn cảnh nào.''\\[4pt]
\textit{(He wrote in the immortal book \textit{Man's Search for Meaning}: `Everything can be taken from a person --- but one thing: the last human freedom is to choose one's attitude toward any given circumstance.')}

\end{truyen}
\begin{baihoc}
  Khi mọi thứ bên ngoài bị tước đoạt, thứ còn lại --- thái độ và ý nghĩa sống --- chính là điều làm nên con người thật sự của bạn.\\[4pt]
  \textit{(When everything outside is stripped away, what remains --- attitude and meaning --- is what truly defines who you are.)}
\end{baihoc}
